\begin{enumerate}
\item $a=-1$, $b=1$, $w(x)=1/\sqrt{1-x^2}$. This leads to the Chebyshev polynomials
\begin{equation}\label{eq:chebyshev-polynomial-cosine-formula}
T_n(x)=\dfrac1{2^{n-1}}\cos(n\arccos x).
\end{equation}

\item $a=-\infty$, $b=\infty$, $w(x)=\exp(-x^2)$. This leads to the Hermite polynomials
\begin{equation}\label{eq:hermite-polynomial-rodrigues-formula}
H_n(x)=(-1)^n\exp(x^2)\dfrac{\mathrm{d}^n}{\mathrm{d}x^n}\exp(-x^2).
\end{equation}

\item $a=0$, $b=\infty$, $w(x)=\exp(-x)$. This leads to the Laguerre polynomials
\begin{equation}\label{eq:laguerre-polynomial-rodrigues-formula}
L_n(x)=e^x\dfrac{\mathrm{d}^n}{\mathrm{d}x^n}(x^ne^{-x}).
\end{equation}
\end{enumerate}

There are various other orthogonal polynomials with specific names, e.g., Jacobi and Gegenbauer polynomials. We leave it to the reader to verify the orthogonality of the Chebyshev, Hermite and Laguerre polynomials; see Problem~\ref{ex:classical-polynomials-orthogonality-normalization}. There are numerous facts known about orthogonal polynomials, e.g., formulas for their coefficients, ``generating functions,'' differential equations which orthogonal polynomials satisfy, recursion relations and relationships to various special functions. We shall not discuss these issues here and instead refer to the literature.

We have yet to address the completeness of the polynomials introduced above. For the Chebyshev polynomials, this is clear from the Weierstrass approximation theorem, since uniform convergence implies convergence in $L_w^2(-1,1)$. For the Hermite and Laguerre polynomials, however, we are dealing with infinite intervals, and we need a somewhat different argument. We first consider the Laguerre polynomials. We have the identity
\begin{equation}\label{eq:laguerre-polynomial-generating-function}
\sum_{n=0}^\infty\dfrac{L_n(x)}{n!}t^n=g(x,t):=\dfrac1{1-t}\exp\left(-\dfrac{xt}{1-t}\right);
\end{equation}
see Problem~\ref{ex:laguerre-generating-function}. Since the Laguerre polynomials are ``generated'' by Taylor expansion of the right-hand side of \eqref{eq:laguerre-polynomial-generating-function}, this expression is referred to as a generating function. We note that the convergence radius of the Taylor series is 1. An explicit calculation shows that
\begin{equation}\label{eq:laguerre-generating-function-lp-two-remainder}
\int_0^\infty e^{-x}\left(g(x,t)-\sum_{n=0}^N\dfrac{L_n(x)}{n!}t^n\right)^2\,\mathrm{d}x=\dfrac1{1-t^2}-\sum_{n=0}^N t^{2n}
\end{equation}


(see Problem~\ref{ex:laguerre-hermite-completeness-details}); hence the series in \eqref{eq:laguerre-polynomial-generating-function} converges also in $L_w^2(0,\infty)$ if $|t|<1$. It follows that any function $e^{-\alpha x}$, $\alpha>-1/2$, can be approximated in $L_w^2(0,\infty)$ by Laguerre polynomials. However, if $f\in L_w^2$ is orthogonal to $e^{-\alpha x}$ for every $\alpha$, then the Laplace transform of $f$ is zero, and therefore $f$ is zero. Hence linear combinations of exponentials $e^{-\alpha x}$ are dense in $L_w^2(0,\infty)$.

For the Hermite polynomials, we have the identity
\begin{equation}\label{eq:hermite-polynomial-generating-function}
e^{-t^2+2tx}=\sum_{n=0}^\infty\dfrac{H_n(x)}{n!}t^n
\end{equation}
corresponding to \eqref{eq:laguerre-polynomial-generating-function} and an analogous argument applies. In this case, the convergence radius of the Taylor series is infinite, and one needs to show that linear combinations of the functions $e^{2tx}$, $t\in\mathbb{C}$, are dense in $L_w^2(-\infty,\infty)$. First observe that $\mathcal{D}(\mathbb{R})$ is dense. Every function $\phi$ in $\mathcal{D}(\mathbb{R})$ can be represented by a convergent Fourier integral:
\begin{equation}\label{eq:hermite-completeness-fourier-inversion}
\phi(x)=\dfrac1{\sqrt{2\pi}}\int_{-\infty}^\infty e^{i\xi x}\widehat\phi(\xi)\,\mathrm{d}\xi,
\end{equation}
and the integral can be approximated by Riemann sums
\begin{equation}\label{eq:hermite-completeness-fourier-riemann-sums}
\phi(x)\sim\dfrac1{\sqrt{2\pi}}\sum_{k=1}^n e^{i\xi_kx}\widehat\phi(\xi_k)(\xi_k-\xi_{k-1}).
\end{equation}
It is easy to see that the discrete sums converge to the integral in the sense of convergence in $L_w^2(-\infty,\infty)$.
